% RaspiMIDIHub manual -- LaTeX header includes
%
% Pulled in via `--include-in-header=templates/header.tex` from the
% `make manual` target. Keep LaTeX-specific tweaks here so the YAML
% metadata stays declarative.

% ------------------------------------------------------------------
% Brand palette. Sourced from website/style.css so the manual and
% the marketing site stay visually aligned.
% ------------------------------------------------------------------
\definecolor{darkBg}{HTML}{0d0d1a}            % cover background
\definecolor{darkBgLight}{HTML}{1a1a2e}       % cover gradient stop
\definecolor{accentColour}{HTML}{00d4aa}      % brand turquoise
\definecolor{accentCoral}{HTML}{e94560}       % brand coral
\definecolor{accentSoft}{HTML}{D9F7F0}        % pale turquoise tint -- table headers
\definecolor{accentCoralSoft}{HTML}{FCE3E8}   % pale coral tint -- warning blocks
\definecolor{lightText}{HTML}{e0e0e0}         % cover body text
\definecolor{mutedText}{HTML}{8892a4}         % cover subtitle / cover author
\definecolor{codeTint}{HTML}{F2F4F7}          % inline code background

% Manual version shown on the cover.
%
% Source of truth: `Makefile`'s `VERSION` variable. `make manual`
% generates `templates/version.tex` from that and pandoc loads it
% via `--include-in-header` BEFORE this file, so `\manualversion`
% is already defined when we get here. These `\providecommand`s
% are the fallback for someone running pandoc directly without
% going through the Makefile -- the placeholder makes it obvious
% the build went off-path.
\providecommand{\manualversion}{(unset -- build via `make manual`)}
\providecommand{\manualdate}{}

% ------------------------------------------------------------------
% Branded cover. Pandoc emits \maketitle; we override it with a dark
% splash that mirrors the website hero. The dark background applies
% only to the title page via eso-pic's starred AddToShipoutPictureBG
% variant. Subsequent pages return to white.
% ------------------------------------------------------------------
\usepackage{eso-pic}
\usepackage{tikz}

\makeatletter
\renewcommand{\maketitle}{%
  \begin{titlepage}%
    \thispagestyle{empty}%
    \AddToShipoutPictureBG*{%
      \AtPageLowerLeft{%
        \begin{tikzpicture}[remember picture, overlay]
          \fill[darkBg] (current page.south west) rectangle (current page.north east);
          \fill[accentColour] (current page.north west) rectangle ([yshift=-8pt]current page.north east);
          \fill[accentCoral] ([yshift=8pt]current page.south west) rectangle (current page.south east);
        \end{tikzpicture}%
      }%
    }%
    \begingroup
    \color{lightText}\sffamily%
    \vspace*{5cm}%
    \begin{center}%
      % Brand mark with the coral "MIDI" middle, exactly like the
      % website hero.
      {\fontsize{56pt}{60pt}\bfseries
        \textcolor{lightText}{Raspi}%
        \textcolor{accentCoral}{MIDI}%
        \textcolor{lightText}{Hub}\par}%
      \vspace{1.2cm}%
      % Subtitle pill -- turquoise rounded background, dark text.
      {\tikz\node[fill=accentColour,
                  text=darkBg,
                  rounded corners=12pt,
                  inner xsep=18pt, inner ysep=8pt,
                  font=\sffamily\bfseries\large]
              {USER MANUAL \enspace$\cdot$\enspace VERSION~\manualversion};%
      \par}%
    \end{center}%
    \vfill
    \begin{center}%
      {\color{mutedText}\large\@author\par}
      \vspace{0.4cm}%
      {\color{mutedText}\large\@date\par}
    \end{center}%
    \vspace*{2.5cm}%
    \endgroup
  \end{titlepage}%
}
\makeatother

% ------------------------------------------------------------------
% Chapter headings: chapter tag in coral, title big and black, with
% a turquoise rule underneath. Mirrors the website's mixed-colour
% accent treatment.
% ------------------------------------------------------------------
% nobottomtitles*: a heading that would land in the last stretch of
% a page (less than \bottomtitlespace remaining) moves to the next
% page instead of sitting orphaned above one or two body lines.
\usepackage[explicit,nobottomtitles*]{titlesec}
\renewcommand{\bottomtitlespace}{0.2\textheight}

\titleformat{\chapter}[display]%
  {\sffamily}%
  {\color{accentCoral}\fontsize{14pt}{16pt}\bfseries
    \MakeUppercase{\chaptertitlename\,\thechapter}}%
  {8pt}%
  {\fontsize{30pt}{34pt}\bfseries\color{black} #1}%
  [\vskip 6pt {\color{accentColour}\hrule height 1.4pt}]%

\titlespacing*{\chapter}{0pt}{-20pt}{30pt}

% Section / subsection numbers in turquoise.
\titleformat{\section}%
  {\sffamily\Large\bfseries}%
  {\color{accentColour}\thesection}%
  {0.8em}%
  {#1}

\titleformat{\subsection}%
  {\sffamily\large\bfseries}%
  {\color{accentColour}\thesubsection}%
  {0.7em}%
  {#1}

\titleformat{\subsubsection}%
  {\sffamily\normalsize\bfseries}%
  {\color{accentColour}\thesubsubsection}%
  {0.6em}%
  {#1}

% ------------------------------------------------------------------
% Table of contents: widen the number boxes so two-digit section /
% subsection numbers (5.10, 8.18, ...) don't collide with the title.
% The defaults in `report.cls` are 2.3em for sections and 3.2em for
% subsections, which fit "1.1" / "1.1.1" only.
% ------------------------------------------------------------------
\makeatletter
\renewcommand*\l@section{\@dottedtocline{1}{1.5em}{3.2em}}
\renewcommand*\l@subsection{\@dottedtocline{2}{4.7em}{3.9em}}
\renewcommand*\l@subsubsection{\@dottedtocline{3}{8.6em}{4.6em}}
\makeatother

% ------------------------------------------------------------------
% Page style: running header with chapter/section names, page number
% in the outer corner. Chapter-opening pages keep the plain style.
% ------------------------------------------------------------------
\usepackage{fancyhdr}
\pagestyle{fancy}
\fancyhf{}
\fancyhead[LE,RO]{\sffamily\thepage}
\fancyhead[LO]{\sffamily\small\nouppercase{\rightmark}}
\fancyhead[RE]{\sffamily\small\nouppercase{\leftmark}}
\renewcommand{\headrulewidth}{0.4pt}
\renewcommand{\footrulewidth}{0pt}

\fancypagestyle{plain}{%
  \fancyhf{}%
  \fancyfoot[C]{\sffamily\thepage}%
  \renewcommand{\headrulewidth}{0pt}%
}

% ------------------------------------------------------------------
% Tables: turquoise rules + tinted header row. Apply \rowcolor in
% the first row of every longtable via a hook. Pandoc emits one
% header row at the top of each longtable; we colour it.
% ------------------------------------------------------------------
\usepackage{longtable}
\usepackage{booktabs}
\usepackage[table]{xcolor}
\arrayrulecolor{accentColour}
\renewcommand{\arraystretch}{1.20}

% Wrap pandoc's \toprule-emitted header in a coloured row by
% redefining \toprule so the first row in each longtable picks up
% the tint without per-table source edits.
\let\origtoprule\toprule
\renewcommand{\toprule}{\origtoprule\rowcolor{accentSoft}}

% ------------------------------------------------------------------
% Inline code: light grey tint background. Wrap pandoc's emitted
% \texttt with a tight colorbox; uses \rule with raisebox so the
% colour rectangle stays within the line.
% ------------------------------------------------------------------
\let\origtexttt\texttt
\renewcommand{\texttt}[1]{%
  \begingroup
    \setlength{\fboxsep}{1.5pt}%
    \colorbox{codeTint}{\origtexttt{\textcolor{black}{#1}}}%
  \endgroup
}

% ------------------------------------------------------------------
% Figure placement. The plugin-reference appendix is wall-to-wall
% portrait phone screenshots (1:2 aspect), which at width=35% end up
% ~10.5 cm tall. Pandoc 3.x emits `\begin{figure}` with no placement
% specifier and metadata fig-pos is ignored. With LaTeX's default
% [tbp] policy, `b` (bottom) gets selected for these figures even
% when the remaining bottom space is too small — they then run past
% the page bottom and clip.
%
% Fix:
%   1. Force figures to TOP or own-page only (\floatplacement{}{tp}).
%      Bottom placement is what was producing the clipping.
%   2. Raise \topfraction so a ~10.5 cm figure (~43% of \textheight)
%      is allowed at the top — the LaTeX default of 0.7 is too tight
%      once the appendix header plus body text are accounted for.
%   3. Same for \floatpagefraction: lower it so a figure that takes
%      40 %+ of a page is still eligible for its own float page
%      instead of being held over indefinitely.
%
% Captions stay sans-serif, accent-coloured label.
% ------------------------------------------------------------------
\usepackage{float}
\floatplacement{figure}{tp}
\renewcommand{\topfraction}{0.9}
\renewcommand{\floatpagefraction}{0.5}
% Keep figures inside the section that references them: [section]
% inserts a \FloatBarrier before every \section, so a deferred
% screenshot cannot drift past the next section heading (it used to
% land under the following plugin's heading in the play-surfaces
% chapter and Appendix A). metadata.yaml's fig-pos comment always
% assumed this barrier; it was missing until now.
\usepackage[section]{placeins}
\usepackage[font=small,labelfont={sf,bf,color=accentColour},
            textfont=sf, justification=centering]{caption}

% ------------------------------------------------------------------
% Code blocks: smaller font, wrap long lines instead of running off
% the page. fvextra is required for the `breakanywhere` option.
% ------------------------------------------------------------------
\usepackage{fvextra}
\DefineVerbatimEnvironment{Highlighting}{Verbatim}{%
  breaklines,%
  breakanywhere,%
  fontsize=\small,%
  commandchars=\\\{\}}
\fvset{breaklines=true, breakanywhere=true, fontsize=\small}

% ------------------------------------------------------------------
% Paragraph spacing: a bit more air between paragraphs, less first-
% line indentation. Reads better at manual length.
% ------------------------------------------------------------------
\usepackage{parskip}

% ------------------------------------------------------------------
% Admonitions. Pandoc emits fenced divs with class .note, .warning,
% or .tip as `\begin{note}` / `\begin{warning}` / `\begin{tip}` --
% Lua filter `templates/admonitions.lua` does the mapping. Here we
% define how each renders: a left rule + soft background + pill
% label inline at the start.
% ------------------------------------------------------------------
\usepackage{mdframed}
\usepackage{tcolorbox}

\newcommand{\admlabel}[2]{%
  \begin{tikzpicture}[baseline=(b.base)]%
    \node[fill=#1, text=white, rounded corners=4pt,
          inner xsep=6pt, inner ysep=2pt,
          font=\sffamily\bfseries\scriptsize] (b) {#2};%
  \end{tikzpicture}\hspace{0.6em}\ignorespaces}

\newmdenv[
  linecolor=accentColour,
  linewidth=3pt,
  topline=false, bottomline=false, rightline=false,
  innerleftmargin=12pt, innerrightmargin=10pt,
  innertopmargin=8pt, innerbottommargin=8pt,
  skipabove=8pt, skipbelow=8pt,
  backgroundcolor=accentSoft,
]{noteframe}

\newmdenv[
  linecolor=accentCoral,
  linewidth=3pt,
  topline=false, bottomline=false, rightline=false,
  innerleftmargin=12pt, innerrightmargin=10pt,
  innertopmargin=8pt, innerbottommargin=8pt,
  skipabove=8pt, skipbelow=8pt,
  backgroundcolor=accentCoralSoft,
]{warningframe}

\newenvironment{note}%
  {\begin{noteframe}\admlabel{accentColour}{NOTE}}%
  {\end{noteframe}}

\newenvironment{warning}%
  {\begin{warningframe}\admlabel{accentCoral}{WARNING}}%
  {\end{warningframe}}

\newenvironment{tip}%
  {\begin{noteframe}\admlabel{accentColour}{TIP}}%
  {\end{noteframe}}
